\subsection{Instruction Payload Ordering}
The instruction (:term:base.terms.header:) is byte 0 for extrashort and short instructions and the first two bytes for medium and
longer instructions. The (:term:base.terms.opcode_space:) may continue after the (:term:base.terms.header:). If a descriptor, immediate, displacement, or
bitmap follows the (:term:base.terms.opcode_space:), each independently determined unit is a separate (:term:base.terms.payload:), and the (:term:base.terms.payload|plural:) occupy
increasing addresses in decode order. An extended EA descriptor and a following displacement are therefore separate
(:term:base.terms.payload|plural:).

Signed displacement and immediate (:term:base.terms.payload|plural:) use two\textquotesingle{}s-complement representation at their declared width before any
sign extension performed by the consuming instruction or effective-address encoding.

\Needspace{1.25in}
\manualtablecaption{Immediate Operand Interpretation}
\begin{manuallongtable}{@{}>{\raggedright\arraybackslash}p{0.95in}>{\raggedright\arraybackslash}p{0.38in}>{\raggedright\arraybackslash}p{0.62in}>{\raggedright\arraybackslash}p{0.62in}>{\raggedright\arraybackslash}p{0.90in}>{\raggedright\arraybackslash}p{1.70in}@{}}
\toprule
\rowcolor{ManualHeaderFill}
\textbf{Operand} & \textbf{Bits} & \textbf{Value} & \textbf{Range} & \textbf{Extension} & \textbf{Applies When}\\
\midrule
\endfirsthead
\multicolumn{6}{@{}l}{\scriptsize\itshape Table \themanualtable\ (continued)}\\
\toprule
\rowcolor{ManualHeaderFill}
\textbf{Operand} & \textbf{Bits} & \textbf{Value} & \textbf{Range} & \textbf{Extension} & \textbf{Applies When}\\
\midrule
\endhead
\texttt{PAIRn} & 3 & identifier & 0..7 & none & PUSHP/POPP canonical register-pair selector\\
\texttt{pt\_level} & 3 & enumeration & 1..5 & none & PTQUERY page-table level selector\\
\texttt{flags\_bitmap} & 4 & bitmap & 0..15 & none & SETF integer FLAGS image\\
\texttt{imm6} & 6 & unsigned & 0..63 & zero-extend & integer encodings with an explicit B/W/L/Q operation size\\
\texttt{imm7} & 7 & unsigned & 0..127 & zero-extend & EXTRACT register-pair encodings\\
\bottomrule
\end{manuallongtable}


PUSHP and POPP interpret their 3-bit operand as a canonical register-pair index, so every encoding
from 0..7 is valid.
PTQUERY uses a 3-bit page-table level; assigned levels are 1..5, while encodings 0, 6, and 7
are (:term:base.terms.reserved:).
SETF interprets its 4-bit bitmap range 0..15 as a complete FLAGS
image: bit 3 writes Z, bit 2 writes N, bit 1 writes C, and bit 0 writes V. Each one writes the corresponding flag to
one, and each zero writes it to zero.

\begin{manuallistedformatdiagram}{SETF Flag Bitmap}
\manualformatrowrange{flag bitmap[3:0]}{3}{0}{%
\manualformatfield{Z}{1}
\manualformatfield{N}{1}
\manualformatfield{C}{1}
\manualformatfield{V}{1}
}
\end{manuallistedformatdiagram}

A 6-bit \texttt{imm6} value in 0..63 is zero-extended to the selected operation size before use when
that size is wider than 6 bits. EXTRACT uses the 7-bit \texttt{imm7} range 0..127 as the
offset into its concatenated register pair.
